\documentclass[10pt,a4paper]{amsart}
\usepackage[margin=19mm]{geometry}
\usepackage{amsmath,amssymb,mathtools}
\usepackage[scr=rsfs,cal=euler]{mathalfa}
\usepackage{xfrac}
\usepackage[unicode,hidelinks,bookmarks=false]{hyperref}
\newcommand{\ie}{i.e.\ }
\newcommand{\cf}{cf.\ }
\newcommand{\resp}{respectively\ }
\newcommand{\Subsection}{\subsection}
\providecommand{\nobreakdash}{\nobreak\hbox{-}}
\setlength{\emergencystretch}{2em}
\multlinegap0pt
\usepackage{bm}
\usepackage{bbm}
\usepackage{dsfont}
\usepackage{xspace}
\usepackage{cancel}
\usepackage{mathtools}
\usepackage[shortlabels]{enumitem}
\usepackage{amsthm}
\makeatletter
\@ifundefined{thm}{\newtheorem{thm}{Theorem}}{}
\@ifundefined{prop}{\newtheorem{prop}[thm]{Proposition}}{}
\makeatother
\DeclareMathOperator{\divv}{div}%%div for divisor
\newcommand{\de}{\delta}
\newcommand{\ka}{\kappa}
\newcommand{\pp}{\mathfrak{p}}
\title[Fibrations by plane projective rational quartic curves in ch]{Fibrations by plane projective rational quartic curves in characteristic two}
\author{Cesar Hilario, Karl-Otto St\"ohr}
\date{Source: arXiv:2409.05464v2 (CC BY 4.0)}
\begin{document}
\maketitle

\noindent\emph{Source and licence.} Fibrations by plane projective rational quartic curves in characteristic two. Version arXiv:2409.05464v2 (CC BY 4.0), distributed under the Creative Commons Attribution 4.0 International licence, as recorded in the arXiv licence metadata for this exact version and shown on the abstract page https://arxiv.org/abs/2409.05464v2. Full rights notice: licenses/arxiv-2409.05464-CC-BY-4.0.txt.\par\smallskip

\noindent\textit{Editorial scope.} Counted unit: the Prop whose source label is \texttt{2024\_08\_04\_19:35} in the article (source0.tex lines 1739--1748, source label \texttt{2024\_08\_04\_19:35}), delivered together with the article's own complete original proof of it (source0.tex lines 1751--1768, one \texttt{proof} environment of 18 lines). No mathematics is written, repaired or re-derived: statement and proof are the article's own text line for line. Required context retained uncounted: 2024\_02\_18\_00:35 (lines 905--926); 2021\_05\_21\_18:40 (lines 1246--1275); 2021\_05\_21\_18:42 (lines 1249--1265); 2021\_05\_21\_18:43 (lines 1249--1265). Every definition, display and label of those lines is retained unchanged. Omitted from this package and not redistributed: 1--904, 927--1245, 1266--1738, 1749--1750, 1769--3101. Nothing inside a retained excerpt is omitted or altered: every retained body is delivered exactly as the article's lines are written, so no omission receipt of any kind accompanies this package. Citations: the retained excerpts carry no citation command at all, so no bibliography entry of the article is transcribed and no reference is redistributed. Reference adaptation: none was needed and none was made; every reference inside the delivered unit resolves inside it, so no unresolved reference or citation is produced. Printed slips kept literal: no printed word, symbol, hypothesis or number of the counted statement or of its proof was altered. Layout and class: the article's own class is \texttt{amsart} with packages including amsmath, amssymb, amsthm, amscd, mathrsfs, old-arrows, stmaryrd, inputenc, babel, tikz, tikz-cd, graphicx, amsrefs, hyperref; the wrapper is the lane's pinned \texttt{amsart} editorial wrapper at 10pt on A4, which restates the theorem-like environments the retained text needs and adds the article's own macro definitions that the retained text needs (\texttt{\textbackslash de}, \texttt{\textbackslash divv}, \texttt{\textbackslash ka}, \texttt{\textbackslash pp}), with extra wrapper packages (bm, bbm, dsfont, xspace, cancel, mathtools, enumitem). The delivered numbering is the unit's own. Assets: no retained span contains a figure environment, a graphics-inclusion command, a PDF-page inclusion, a table or a TikZ picture; no image file is copied into this package, the package declares no asset dependency and no image file is redistributed with it. Licence: version arXiv:2409.05464v2 of this article is distributed under the Creative Commons Attribution 4.0 International licence, as recorded in the arXiv licence metadata for this exact version and shown on the abstract page https://arxiv.org/abs/2409.05464v2. A full rights notice is delivered at licenses/arxiv-2409.05464-CC-BY-4.0.txt. Characters: every retained excerpt is pure ASCII, so the delivered engine, XeLaTeX with its default Latin Modern text font, reports no missing character.
\begin{thm}\label{2024_02_18_00:35}
A one-dimensional separable function field $F|K$ of characteristic $p=2$ and genus $g = 3$ is geometrically rational and admits a prime $\pp$ such that $\de(\pp) = 3$, $\de(\pp_1) = 1$, $\pp_2$ is rational, and such that $\pp$ is not a canonical divisor, if and only if 
$F=K(x,z,y)$ is generated by three functions $x$, $z$, and $y$ that can be put into the following normal form
%\todoo{change notation $A\gg a$, $B \gg b$?}
\begin{align*}
z^2 &= (c_0 + c_1 x + x^2) (c_0 A_2 + c_1^{-1} + c_1 A_2 x + A_2 x^2), \\
y^2 &= (c_0 + c_1 x + x^2) (B_0 + B_1 x + z),
\end{align*}
%\[ z^2 = c(x) A(x), \quad y^2 = c(x) (B(x) + z),  \]
%where the polynomials $c(x)$, $A(x)$ and $B(x)$ are given by
%\begin{align*}
%    c(x) &= c_0 + c_1 x + x^2,  \\ 
%    A(x) &= (c_0 A_2 + c_1^{-1}) + c_1 A_2 x + A_2 x^2,  \\ 
%    B(x) &= B_0 + B_1 x,
%\end{align*}
%and the constants $c_0,c_1,A_2,B_0,B_1\in K$ satisfy the conditions $c_1\neq 0$ and $A_2\notin K^2$.
where $c_0,c_1,A_2,B_0,B_1\in K$ are constants satisfying the conditions $c_1\neq 0$ and $A_2\notin K^2$.
The singular prime $\pp$ is the only pole of the function $x$.
It has 
degrees $\deg(\pp)=4$, $\deg(\pp_1)=2$, $\deg(\pp_2)=1$, and
residue fields $\ka(\pp)=\ka(A_2^{1/4})$, $\ka(\pp_1)=K(A_2^{1/2})$, $\ka(\pp_2)=K$.
\end{thm}

\begin{thm}\label{2021_05_21_18:40}
A one-dimensional separable function field $F|K$ of characteristic $p=2$ and genus $g = 3$ is geometrically rational and admits a prime $\pp$ such that $\de(\pp) = 3$, $\de(\pp_1) = 1$, $\pp_2$ is non-rational, and such that $\pp$ is not a canonical divisor, if and only if 
$F|K$ can be put into one of the following normal forms
\begin{enumerate}[\upshape (a)]
    \setcounter{enumi}{1}
    \item \label{2021_05_21_18:42}
    $F|K = K(x,w,z,y)|K$, where
    \[ w^2 = x + a_2 x^2, \quad z^2 = b_0 + b_2 x^2 + w, \quad y^2 = xz, \]
    and $a_2,b_0,b_2 \in K$ are constants satisfying $a_2 \notin K^2$.
    
    \item \label{2021_05_21_18:43}
    $F|K = K(x,w,z,y)|K$, where
    \[ w^2 = a_ 0 + x + a_2 x^2, \quad z^2 = b_1 w^2 + w, \quad y^2 = (c_3 + c_4 x + z) w, \]
    and $a_0,a_2,b_1,c_3,c_4 \in K$ are constants satisfying $a_2,b_1 \notin K^2$.

    \item \label{2021_05_21_18:45}
    $F|K = K(x,z,y)|K$, where    
    \[ z^4 = a_0 + x + a_2 x^2, \quad y^2 = c_0 + z + c_2 z^2, \]
    and $a_0,a_2,c_0,c_2 \in K$ are constants satisfying $a_2\notin K^2$ and $c_2 \in K^2 a_2$.
\end{enumerate}
In each case the singular prime $\pp$ is the only pole of the function $x$.
It has residue fields $\ka(\pp_3) = K$, $\ka(\pp_2) = K(a_2^{1/2})$,
and
\begin{enumerate}[\upshape (a)]
    \setcounter{enumi}{1}
    \item $\ka({\pp_1}) = K(a_2^{1/2}, b_2^{1/2})$, $\ka(\pp) = K(a_2^{1/2}, b_2^{1/4})$,
    \item $\ka({\pp_1}) = K(a_2^{1/2}, b_1^{1/2})$, $\ka(\pp) = K(a_2^{1/2}, b_1^{1/2}, (c_4 a_2^{1/2} + a_2 b_1^{1/2})^{1/2})$,
    \item $\ka(\pp_1) = \ka(\pp) = K(a_2^{1/2})$.
\end{enumerate}
\end{thm}

\begin{prop}\label{2024_08_04_19:35}
Let $F|K$ be a function field as in Theorem~\ref{2024_02_18_00:35}, Theorem~\ref{2021_05_21_18:40}~\ref{2021_05_21_18:42} or Theorem~\ref{2021_05_21_18:40}~\ref{2021_05_21_18:43}.
Then a basis of the vector space of holomorphic differentials of $F|K$ is given by
\begin{enumerate}[\upshape (a)]
    \item $dy, \, \breve z \, dy, \, \breve y \, dy$, where $\breve z:= z \cdot c(x)^{-1}$, $\breve y := y \cdot c(x)^{-1}$, $c(x) := c_0 + c_1 x + x^2$, in Theorem~\ref{2024_02_18_00:35};
    \item $dy$, $\breve w \, dy$, $\breve y \, dy$, where $\breve w := wx^{-1}$, $\breve y := y x^{-1}$, in Theorem~\ref{2021_05_21_18:40}~\ref{2021_05_21_18:42};
    \item $dy$, $\breve z \, dy$, $\breve y \, dy$, where $\breve z := zw^{-1}$, $\breve y := yw^{-1}$, in Theorem~\ref{2021_05_21_18:40}~\ref{2021_05_21_18:43}.
\end{enumerate}
In each case the first two differentials form a basis of the vector space of exact holomorphic differentials of $F|K$.
\end{prop}

\begin{proof}
Suppose $F|K$ is given as in Theorem~\ref{2024_02_18_00:35}.
Since the functions $\breve y$ and $\breve z$ fulfill the relations $\breve z^2 = A_2 + c_1^{-1} \cdot c(x)^{-1}$ and $\breve y^2 = (B_0 + B_1 x) \cdot c(x)^{-1} + \breve z$,
their pole divisors satisfy $\mathrm{div}_\infty(\breve z) = \divv_0(c(x))^{1/2}$ and $\mathrm{div}_\infty(\breve y) \leq \divv_0(c(x))^{1/2}$.
By the preceding lemma, this implies that the three differentials are holomorphic. As they are linearly independent and as their number is equal to the dimension $g=3$ of the vector space of holomorphic differentials, they form a basis of this vector space.

For $F|K$ as in Theorem~\ref{2021_05_21_18:40}, \ref{2021_05_21_18:42} or~\ref{2021_05_21_18:43}, we argue analogously.
If $F|K$ is of type~\ref{2021_05_21_18:42} (resp. type~\ref{2021_05_21_18:43}), then the relations
$\breve w^2 = a_2 + x^{-1}$ and $\breve y^4 = b_2 + b_0 x^{-2} + \breve w x^{-1}$
(resp. $\breve z^2 = b_1 + w^{-1}$, $w^2 = a_0 + x + a_2 x^2$ and $\breve y^2 = c_3 w^{-1} + c_4 x w^{-1} + \breve z$)
show that
$\divv_\infty(\breve w) = \divv_0(x)^{1/2}$ and $\divv_\infty(\breve y) \leq \divv_0(x)^{1/2}$
(resp. $\divv_\infty (\breve z) = \divv_0(w)^{1/2}$ and $\divv_\infty(\breve y)\leq \divv_0(w)^{1/2}$),
that is, the three differentials are holomorphic and form a basis of the vector space of holomorphic differentials of $F|K$.

It remains to note that as $F=F_1 \oplus F_1 y$, the space of exact differentials of $F|K$ is equal to $F_1 \, dy$.
Hence in each case the first two elements of the basis are exact, while the third element is not exact.
\end{proof}

\end{document}
